\section{Norm bounds for restoring operators}
\label{sec:restoring_norm}

Let $X,Y$ be finite-dimensional complex Hilbert spaces, fix a full
orthonormal basis $\{|x\rangle\}_{x\in\mathsf X}$ of $X$, and let
$E\subseteq\mathsf X$. Use the operators $R_0,B_0,B_Y$ of
Definition~\ref{def:restoring_operators}, with the same tensor orders.
All vector norms are Euclidean norms, and all operator norms are the
induced norms. For the weighted-sum inequalities alone, the numbers
$p_x$ may be arbitrary real numbers, with $0^{-1}=0$; positivity on $E$
will be imposed for the norm bounds. These estimates give the first
restoring-operator estimate in the amplification argument
of~\cite{OpenAI2026AreaLaw}.

\begin{theorem}[Selected blocks below the full partial trace]
    \label{thm:restoring_selected_trace}
    \lean{Entropy.restoringGram_le_compressed_partialTrace}
    \leanok
    \uses{def:restoring_operators, def:partial_trace}
    Suppose $Q\ge0$, $T=T^\dagger$, and $c_X\ge0$.
    If $p_x^{-1}\le c_X$ for every $x\in E$, then
    \begin{align}
        B_Y\le c_X T(\operatorname{tr}_X Q)T.
        \label{eq:restoring_selected_trace}
    \end{align}
\end{theorem}
\begin{proof}\leanok
    \uses{def:restoring_operators}
    Put $Q_x=\langle x|Q|x\rangle$. Each $Q_x$ and $TQ_xT$ is
    positive semidefinite. Therefore
    \begin{align*}
        B_Y
        &\le c_X T\left(\sum_{x\in E}Q_x\right)T
        \le c_X T\left(\sum_{x\in\mathsf X}Q_x\right)T
        =c_X T(\operatorname{tr}_X Q)T.
    \end{align*}
    The second inequality follows by summing the positive omitted blocks
    and taking congruence by $T$.
\end{proof}

\begin{theorem}[Spectral bound for the restoring Gram factor]
    \label{thm:restoring_gram_bound}
    \lean{Entropy.restoringGram_le}
    \leanok
    \uses{def:restoring_operators, def:partial_trace}
    Let $\rho$ be an operator on $X\otimes Y$, and set
    $\rho_Y=\operatorname{tr}_X\rho$. Suppose $Q$ and $T$ are
    orthogonal projections and $c_X,c_Q,c_Y\ge0$. Assume the following
    inequalities, with the scalar inequality required for every $x\in E$:
    \begin{align}
        p_x^{-1}&\le c_X,\label{eq:restoring_spectral_bounds}\\
        Q&\le c_Q\rho,\notag\\
        T\rho_YT&\le c_YT.\notag
    \end{align}
    Then $B_Y\le c_Xc_Qc_Y I_Y$.
\end{theorem}
\begin{proof}\leanok
    \uses{thm:restoring_selected_trace}
    Positivity of the partial trace gives
    $\operatorname{tr}_XQ\le c_Q\rho_Y$.
    Apply~(\ref{eq:restoring_selected_trace}), preserve order under
    congruence by $T$, and use $T\le I_Y$ to obtain
    \begin{align*}
        B_Y
        &\le c_X T(\operatorname{tr}_XQ)T
        \le c_Xc_Q T\rho_YT
        \le c_Xc_Qc_Y T
        \le c_Xc_Qc_Y I_Y.
    \end{align*}
\end{proof}

\begin{theorem}[Outer product of a unit blank]
    \label{thm:restoring_unit_blank}
    \lean{Entropy.blankOuterProduct_isStarProjection}
    \leanok
    \uses{def:restoring_operators}
    If $b\in X$ and $\|b\|=1$, then $P_b=|b\rangle\langle b|$
    is an orthogonal projection.
\end{theorem}
\begin{proof}\leanok
    \uses{def:restoring_operators}
    The outer product is positive semidefinite and hence Hermitian.
    Its square is $P_b^2=\langle b,b\rangle P_b=\|b\|^2P_b=P_b$.
\end{proof}

\begin{theorem}[Norm bounds from the spectral constants]
    \label{thm:restoring_operator_norm}
    \lean{Entropy.restoringOperator_norm_le, Entropy.columnOperator_norm_le}
    \leanok
    \uses{def:restoring_operators, def:partial_trace}
    Let $\rho$ be an operator on $X\otimes Y$, and set
    $\rho_Y=\operatorname{tr}_X\rho$. Suppose $Q,T$ are orthogonal
    projections, $c_X,c_Q,c_Y\ge0$, and
    (\ref{eq:restoring_spectral_bounds}) holds. Assume $p_x>0$ for
    every $x\in E$.
    If $\|0_s\|=\|0_X\|=1$, then $\|R_0\|\le\sqrt{c_Xc_Qc_Y}$.
    Independently, if $\|0_s\|=\|0_e\|=1$, then
    $\|B_0\|\le\sqrt{c_Xc_Qc_Y}$.
    Neither bound requires positivity of the unselected $p_x$ or
    commutation between $Q$ and $I_X\otimes T$.
\end{theorem}
\begin{proof}\leanok
    \uses{thm:restoring_gram_bound, thm:restoring_unit_blank, thm:restoring_grams}
    Set $C=c_Xc_Qc_Y$. Theorem~\ref{thm:restoring_grams} gives $B_Y\ge0$,
    and Theorem~\ref{thm:restoring_gram_bound} gives $B_Y\le CI_Y$.
    Each normalized blank gives $0\le P_b\le I_X$ by
    Theorem~\ref{thm:restoring_unit_blank}.
    For positive $A,B$ satisfying $A\le I$ and $B\le cI$, positivity
    of tensor products yields $A\otimes B\le I\otimes B\le cI$.
    Apply this inequality to the tensor factors in both Gram formulas
    (\ref{eq:restore_rgram}) and~(\ref{eq:restore_bgram}). Thus
    \begin{align*}
        R_0^\dagger R_0&\le CI_{s\otimes(X\otimes Y)},\\
        B_0^\dagger B_0&\le CI_{(s\otimes e)\otimes(X\otimes Y)}.
    \end{align*}
    For either operator $A$, positivity and this order bound give
    $\|A\|^2=\|A^\dagger A\|\le C$. Taking the nonnegative square
    root proves the two conclusions.
\end{proof}

\begin{theorem}[Exponential specialization]
    \label{thm:restoring_norm_exp}
    \lean{Entropy.restoringOperators_norm_le_exp}
    \leanok
    \uses{def:restoring_operators, def:partial_trace}
    Let $\rho$ be an operator on $X\otimes Y$, set
    $\rho_Y=\operatorname{tr}_X\rho$, and suppose $Q,T$ are
    orthogonal projections. Assume $p_x>0$ for every $x\in E$ and
    $\|0_s\|=\|0_X\|=\|0_e\|=1$.
    Let $S_X,S_{XY},S_Y,w,\mathcal I$ be real numbers satisfying
    $\mathcal I=S_X+S_{XY}-S_Y$. Assume the following inequalities,
    with the scalar inequality required for every $x\in E$:
    \begin{align}
        p_x^{-1}&\le e^{S_X+w},\label{eq:restoring_exp_hypotheses}\\
        Q&\le e^{S_{XY}+w}\rho,\notag\\
        T\rho_YT&\le e^{-S_Y+w}T.\notag
    \end{align}
    Then
    \begin{align}
        \|R_0\|&\le\exp\bigl((\mathcal I+3w)/2\bigr),
        \label{eq:restoring_norm_exp}\\
        \|B_0\|&\le\exp\bigl((\mathcal I+3w)/2\bigr).\notag
    \end{align}
    In the spectral application, take
    $S_X=S(X)$, $S_{XY}=S(XY)$, and $S_Y=S(Y)$.
\end{theorem}
\begin{proof}\leanok
    \uses{thm:restoring_operator_norm}
    Apply Theorem~\ref{thm:restoring_operator_norm} with
    $c_X=e^{S_X+w}$, $c_Q=e^{S_{XY}+w}$, and $c_Y=e^{-S_Y+w}$.
    All three constants are positive, and their square-root product is
    \begin{align*}
        \sqrt{e^{S_X+w}e^{S_{XY}+w}e^{-S_Y+w}}
        &=\exp\bigl((S_X+S_{XY}-S_Y+3w)/2\bigr)
        =\exp\bigl((\mathcal I+3w)/2\bigr).
    \end{align*}
    This is the scalar specialization of the spectral-constant estimate.
\end{proof}

% Primary source: adc7f1241b42e322a6451854ab7e4b4c146bf78a,
% preprints/A-two-dimensional-area-law-from-a-global-spectral-gap-September-24-2026/
% build/sections/09-amplification.tex:374--397.
% The scalar identity defining I is a hypothesis here. No claim about the
% complete amplification argument or a purity-to-mutual-information identity
% is attached to this leaf.
